\documentclass[12pt]{article}
\usepackage[obeyspaces]{url}
\usepackage{amsmath,amssymb}
\usepackage{fullpage}
\usepackage{hyperref}
\usepackage{mathtools}
\usepackage{seqsplit}
\usepackage{tikz}
\usepackage{MnSymbol,wasysym}
\usepackage{pdflscape}
\usepackage{lscape}
\usepackage{newpxtext}
\usepackage{booktabs}

\setlength{\parindent}{0pt}
\setlength{\parskip}{2ex}
\thispagestyle{empty}

\hypersetup{
    colorlinks=true,
    linkcolor=blue,
    filecolor=magenta,      
    urlcolor=cyan,
    pdftitle={Overleaf Example},
    pdfpagemode=FullScreen,
    }

\newcommand{\N}{\ensuremath{\mathbb{N}}}
\newcommand{\Q}{\ensuremath{\mathbb{Q}}}
\newcommand{\R}{\ensuremath{\mathbb{R}}}
\newcommand{\Z}{\ensuremath{\mathbb{Z}}}

\begin{document}

\includegraphics{CSUSig}

\begin{center}
    {\large\textbf{MATH 101 - Section 002 - Final Poster Project}}
\end{center}
\textbf{\underline{Instructions:}} Working in groups of 2 or 3, you will prepare a poster about a topic related to math in the social sciences. You will present your poster at our Math 101 Poster Session on \textcolor{red}{\textbf{Tuesday, May 13th, 7:30am - 9:30am}, in our usual classroom (Nutrien 140)}. This is our final exam date and is set by the university and cannot be changed.  You must be at the poster session to receive full credit on your project. During the poster session, I will walk around and ask you to explain your poster to me. You will also walk around, talk to your peers about their posters, and submit reflections on three 3 other posters. Key dates and grading scheme is as follows:

\begin{itemize}
    \item \textcolor{red}{April 16th}, 11:59pm: Submit group plans on Canvas (Tell me your group or tell me you want me to match you with someone) (1 point)
    \item \textcolor{red}{April 23rd}, 11:59pm: Submit project topic on Canvas (2 points)
    \item \textcolor{red}{May 13th}, 7:30am: Poster session in \textcolor{red}{room}
        \begin{itemize}
            \item 7:30am: Submit poster on Canvas \textbf{before} the poster session (1 point)
            \item 7:30am - 9:30am: Participate in poster session (14 points for clarity, correctness, background, references, and verbal explanation)
             \item 9:30am: Submit reflections on paper (3 points)
             \end{itemize}
\end{itemize}

\textbf{\underline{Format:}} Your poster can be printed or drawn or a combination of the two. You can be creative with shape and materials. Your poster should include your names, poster title, and references (where you got your information). Before the poster session, you will upload a picture of your poster as a \textbf{.pdf} or \textbf{.jpg} file to Canvas. 

\textbf{\underline{Topic:}}
Your poster topic needs to go beyond the scope of the class, so your topic need to be something we didn't cover in class.  A few ideas are: \\ You can take any topic that we have covered and find \textbf{real-world} examples. You can do a poster on topics in the book we \textbf{haven't} covered. You can \textbf{extend} a topic we learned about in class.  You can find math in something that \textbf{interests you} (music, sports, art, biology, weather, finances, crime, AI, space travel, etc.). Google ``math in...'' to see what comes up. Questions to think about are:
\begin{itemize}
    \item What problems were people attempting to solve?
    \item What background or math do we need to know to understand this topic?
    \item How did society (people, current events, technology, etc.) influence this topic?
\end{itemize}

You will submit a project proposal to me on \textcolor{red}{April 23rd} with your chosen topic, and I will approve it or offer suggestions. If you tell me your interests, I am happy to offer topic ideas!\par

\textbf{\underline{Poster Advice}}
The principles for a poster are similar to those of a presentation, meaning you should prioritize pictures and short descriptions. Unlike a presentation, however, a poster needs to be understood without you there, so some of your descriptions can be a bit longer than what you would use for a presentation.

Here is some advice taken from Anya Michaelsen's excellent \href{https://www.ams.org/journals/notices/202002/rnoti-p189.pdf}{article} on the topic:
\begin{itemize}
    \item \textbf{Structure your Poster Effectively:} Organize your poster into clear and defined sections, using boxes and titles. Follow a logical flow, typically from top to bottom and left to right. Consider sections like Introduction, Background, Results, and Discussion.

    \item \textbf{Use White Space:} Avoid overwhelming your poster with dense text. Keep text concise and use bullet points or lists. Utilize headings, subheadings, and space between elements for clarity.

    \item \textbf{Incorporate Visual Elements:} Use color, bold, italics, and other formatting to highlight key points. Include figures, diagrams, or pictures to break up text and enhance understanding. Consider color coding to differentiate important concepts.

    \item \textbf{Tailor Content to Your Audience:} Determine your intended audience's background and knowledge level. Adjust the level of detail in your background section accordingly. Ensure your poster can stand alone if necessary or have an accompanying one-minute summary.
\end{itemize}

% While not mentioned in Anya's article, you should also cite your sources. You can cite in any format you wish.

\underline{\textbf{Templates:}} You don't have to design your poster from scratch. You can, and probably should, use a template to get started. Spend time thinking about how you want your poster to look before choosing a template. Below I linked some resources for templates, but there are many other places on the internet to get templates.
\begin{itemize}
    \item University of Minnesota: \url{https://libguides.umn.edu/posters/googleslides}
    \item Western Carolina University: \url{https://researchguides.wcu.edu/researchposter/layout}
    \item The Craft of Scientific Presentations: \url{https://www.craftofscientificposters.com/templates.html}
    \item PosterPresentations.com: \url{https://www.posterpresentations.com/how-to-add-our-poster-templates-to-google-slides.html}
\end{itemize}

\underline{\textbf{Printing:}} Every college within the university handles printing posters separately. You may be able to get your poster printed at no cost. Follow the directions for your college here: 

\url{https://tilt.colostate.edu/oura/poster-printing-instructions-for-curc/}

\underline{\textbf{Presenting the poster:}} In addition to making the poster, you should prepare a 2--3 minute presentation, explaining the main points of the poster. Each of your group members should have a speaking part during your presentation. You will present your poster to the instructors and also your peers.

\underline{\textbf{Examples:}} You will find examples on Canvas. Take note of things you like and don't like.

\textbf{\underline{Grading Rubric:}} With a total of 20 points, please see next page.
\begin{table*}[h!]
    \centering
    \begin{tabular}{p{0.4\linewidth}p{0.4\linewidth}p{0.05\linewidth}}\toprule
        \textbf{Criteria} & \textbf{Marks} & \textbf{Pts.}\\ \midrule
        {\textbf{Clarity of poster}:\par Easily legible, clear flow of ideas, has headings, images, title, etc.}& 2 (meets expectations)\par 0.67 (below expectations)\par 0 (missing)& 2\\\midrule
        \textbf{Math}:\par Math involved is correct,  explained and meaningful to the topic& 3 (meets expectations)\par 2 (near expectations)\par 1 (below expectations)\par 0 (missing) & 3\\\midrule
        \textbf{Background}:\par Math or other content needed to understand the rest of the poster is explained & 3 (meets expectations)\par 2 (near expectations)\par 1 (below expectations)\par 0 (missing) & 3\\\midrule
        \textbf{References}:\par References are included on poster, can be cited in any format you wish. & 2 (meets expectations)\par 1 (below expectations)\par 0 (missing)& 2\\\midrule
        \textbf{Explanation}\par Group is able to explain and answer questions while people walk around and see posters.&3 (meets expectations)\par 2 (near expectations)\par 1 (below expectations)\par 0 (missing) & 3\\\midrule
        \textbf{Peer poster reflections}:\par Write reflections on $3$ posters you saw at the poster session.& 3 (meets expectations)\par 2 (near expectations)\par 1 (below expectations)\par 0 (missing) & 3\\\midrule
        \textbf{Timely submission of group plans}:\par Submitted the list of group members or asked to be grouped by the due date. & 1 (did)\par 0 (didn't)& 1\\\midrule
        \textbf{Timely submission of proposal}:\par Submitted the project proposal by the due date with a small explanation on the main idea& 2 (meets expectations)\par 1 (below expectations)\par 0 (Missing)& 2\\\midrule
        \textbf{Timely submission of poster}:\par Submitted the poster by the due date and in the correct format.& 1 (did)\par 0 (didn't)& 1\\
   \bottomrule
    \end{tabular}
    \end{table*}
\end{document}